% GENERATED FILE. Edit canonical lessons, then run npm run generate:books.
% canonical-source-hash: fnv1a64:c84107b44ba7e209
% canonical-lessons: ES-C752-ajedrez, ES-C752-rompecabezas, ES-C752-baraja, ES-C752-magia, ES-C752-apostar

\chapter{Juegos de mesa}
\label{ch:es-juegos-de-mesa}
\hlchaptermodality{752}

\begin{chapteropening}
\textbf{By the end of this chapter:} \emph{I can talk about board games, cards, puzzles, magic tricks and betting.}

Complete the last lesson of chapter 752: I can talk about board games, cards, puzzles, magic tricks and betting.
\end{chapteropening}

\section[el ajedrez]{el ajedrez — chess}

\label{lesson:ES-C752-ajedrez}



\begin{quote}
Before the new one: say the Spanish for the collection, then the Spanish for watercolour.
\end{quote}

\subsection*{You'll want to know: el ajedrez}
\textbf{el ajedrez} — \textquotedblleft{}chess\textquotedblright{}.

\begin{quote}
Juego al ajedrez con mi abuelo.
\end{quote}

\begin{quote}
\emph{I play chess with my grandfather.}
\end{quote}

\emph{Jugar al ajedrez} (in Latin America also \emph{jugar ajedrez}).

\begin{cousinweb}
\emph{ajedrez} was borrowed from Andalusian Arabic \emph{aš-šaṭranj} (EtymDB), which came from Persian and, further back, from Sanskrit \emph{caturaṅga} \textquotedblleft{}the four-limbed army\textquotedblright{}.
\end{cousinweb}

\begin{grammarlens}[title={What you've built}]
One more word for telling, explaining and describing at B1.
\end{grammarlens}

\subsection*{Your turn}
\begin{itemize}
\raggedright
  \item \emph{Say it:} \emph{el ajedrez}
  \item \emph{Say it:} \emph{el ajedrez}, once more
  \item \emph{Say it:} a sentence of your own with \emph{el ajedrez}
  \item \emph{Recall:} say the Spanish for the collection, then the Spanish for watercolour, then say \emph{el ajedrez} again
\end{itemize}

\begin{tcolorbox}[breakable,colback=teal!4,colframe=teal!35!black,title={Before you move on}]
What does el ajedrez mean? (\textquotedblleft{}Chess\textquotedblright{}.) Say it once more.
\end{tcolorbox}
\section[el rompecabezas]{el rompecabezas — the jigsaw puzzle; the puzzle}

\label{lesson:ES-C752-rompecabezas}



\begin{quote}
Before the new one: say the Spanish for watercolour, then the Spanish for chess.
\end{quote}

\subsection*{You'll want to know: el rompecabezas}
\textbf{el rompecabezas} — \textquotedblleft{}the jigsaw puzzle; the puzzle\textquotedblright{}.

\begin{quote}
Hicimos un rompecabezas de mil piezas.
\end{quote}

\begin{quote}
\emph{We did a thousand-piece jigsaw.}
\end{quote}

Same form in the plural. In Spain also \emph{el puzle}.

\begin{cousinweb}
\emph{rompecabezas} is built on \emph{romper} \textquotedblleft{}to break\textquotedblright{} and \emph{cabeza} \textquotedblleft{}head\textquotedblright{}, from Latin \emph{rumpere} and \emph{caput} (EtymDB).
\end{cousinweb}

\begin{grammarlens}[title={What you've built}]
One more word for telling, explaining and describing at B1.
\end{grammarlens}

\subsection*{Your turn}
\begin{itemize}
\raggedright
  \item \emph{Say it:} \emph{el rompecabezas}
  \item \emph{Say it:} \emph{el rompecabezas}, once more
  \item \emph{Say it:} a sentence of your own with \emph{el rompecabezas}
  \item \emph{Recall:} say the Spanish for watercolour, then the Spanish for chess, then say \emph{el rompecabezas} again
\end{itemize}

\begin{tcolorbox}[breakable,colback=teal!4,colframe=teal!35!black,title={Before you move on}]
What does el rompecabezas mean? (\textquotedblleft{}The jigsaw puzzle; the puzzle\textquotedblright{}.) Say it once more.
\end{tcolorbox}
\section[la baraja]{la baraja — the pack of cards}

\label{lesson:ES-C752-baraja}



\begin{quote}
Before the new one: say the Spanish for chess, then the Spanish for the jigsaw puzzle.
\end{quote}

\subsection*{You'll want to know: la baraja}
\textbf{la baraja} — \textquotedblleft{}the pack of cards\textquotedblright{}.

\begin{quote}
Saca la baraja y jugamos.
\end{quote}

\begin{quote}
\emph{Get the cards out and let us play.}
\end{quote}

\emph{Barajar} — to shuffle. \emph{La baraja española} has 40 or 48 cards.

\begin{cousinweb}
\emph{baraja} is built on \emph{barajar} (EtymDB), which in Old Spanish meant to quarrel, then to mix up and shuffle. Beyond that the origin is not settled.
\end{cousinweb}

\begin{grammarlens}[title={What you've built}]
One more word for telling, explaining and describing at B1.
\end{grammarlens}

\subsection*{Your turn}
\begin{itemize}
\raggedright
  \item \emph{Say it:} \emph{la baraja}
  \item \emph{Say it:} \emph{la baraja}, once more
  \item \emph{Say it:} a sentence of your own with \emph{la baraja}
  \item \emph{Recall:} say the Spanish for chess, then the Spanish for the jigsaw puzzle, then say \emph{la baraja} again
\end{itemize}

\begin{tcolorbox}[breakable,colback=teal!4,colframe=teal!35!black,title={Before you move on}]
What does la baraja mean? (\textquotedblleft{}The pack of cards\textquotedblright{}.) Say it once more.
\end{tcolorbox}
\section[la magia]{la magia — magic}

\label{lesson:ES-C752-magia}



\begin{quote}
Before the new one: say the Spanish for the jigsaw puzzle, then the Spanish for the pack of cards.
\end{quote}

\subsection*{You'll want to know: la magia}
\textbf{la magia} — \textquotedblleft{}magic\textquotedblright{}.

\begin{quote}
Le encantan los trucos de magia.
\end{quote}

\begin{quote}
\emph{He loves magic tricks.}
\end{quote}

\emph{Por arte de magia} — as if by magic.

\begin{cousinweb}
\emph{magia} comes from Late Latin \emph{magīa}, from Greek \emph{mageía} (EtymDB), built on \emph{mágos}, a Persian priest.
\end{cousinweb}

\begin{grammarlens}[title={What you've built}]
One more word for telling, explaining and describing at B1.
\end{grammarlens}

\subsection*{Your turn}
\begin{itemize}
\raggedright
  \item \emph{Say it:} \emph{la magia}
  \item \emph{Say it:} \emph{la magia}, once more
  \item \emph{Say it:} a sentence of your own with \emph{la magia}
  \item \emph{Recall:} say the Spanish for the jigsaw puzzle, then the Spanish for the pack of cards, then say \emph{la magia} again
\end{itemize}

\begin{tcolorbox}[breakable,colback=teal!4,colframe=teal!35!black,title={Before you move on}]
What does la magia mean? (\textquotedblleft{}Magic\textquotedblright{}.) Say it once more.
\end{tcolorbox}
\section[apostar]{apostar — to bet}

\label{lesson:ES-C752-apostar}



\begin{quote}
Before the new one: say the Spanish for the pack of cards, then the Spanish for magic.
\end{quote}

\subsection*{You'll want to know: apostar}
\textbf{apostar} — \textquotedblleft{}to bet\textquotedblright{}.

\begin{quote}
Te apuesto a que gana el Madrid.
\end{quote}

\begin{quote}
\emph{I bet you Madrid win.}
\end{quote}

\emph{Te apuesto a que...} — I bet you that... Stem change: \emph{apuesto}.

\begin{cousinweb}
\emph{apostar} is usually explained as built on \emph{puesto} \textquotedblleft{}put, placed\textquotedblright{}, the participle of \emph{poner}, which comes from Latin \emph{pōnere} \textquotedblleft{}to put\textquotedblright{}: to bet is to put money down.
\end{cousinweb}

\begin{grammarlens}[title={What you've built}]
One more word for telling, explaining and describing at B1.
\end{grammarlens}

\subsection*{Your turn}
\begin{itemize}
\raggedright
  \item \emph{Say it:} \emph{apostar}
  \item \emph{Say it:} \emph{apostar}, once more
  \item \emph{Say it:} a sentence of your own with \emph{apostar}
  \item \emph{Recall:} say the Spanish for the pack of cards, then the Spanish for magic, then say \emph{apostar} again
\end{itemize}

\begin{tcolorbox}[breakable,colback=teal!4,colframe=teal!35!black,title={Before you move on}]
What does apostar mean? (\textquotedblleft{}To bet\textquotedblright{}.) Say it once more.
\end{tcolorbox}
